
\documentclass[11pt,a4paper]{article}

% Cargar configuración desde archivo externo
\input{Config}

% =========================
% DOCUMENTO
% =========================
\begin{document}
\color{textBrown}

% ============================================================
% PORTADA
% ============================================================
\colorbar
\pagestyle{empty}

\vspace{1.5cm}

% Logos (placeholder - reemplazar con logos reales)
\noindent
\begin{minipage}[c]{0.2\textwidth}
    \centering
    % Logo IEEE (diamante)
    \includegraphics[width=0.9\linewidth]{IEEE.png}
\end{minipage}%
\hfill
\begin{minipage}[c]{0.4\textwidth}
    \centering
    % Logo IEEE Aguascalientes Section
    \fbox{\parbox{4cm}{\centering\vspace{0.3cm}IEEE Aguascalientes\\Section\vspace{0.3cm}}}
\end{minipage}%
\hfill
\begin{minipage}[c]{0.25\textwidth}
    \centering
    % Logo Young Professionals
    \fbox{\parbox{2.5cm}{\centering\vspace{0.3cm}IEEE Young\\Professionals\vspace{0.3cm}}}
\end{minipage}

\vspace{3cm}

% Título principal
{\fontsize{36}{42}\selectfont\bfseries\color{titleBrown} Plan estratégico de trabajo 2026}

\vspace{0.8cm}

% Subtítulo
{\large Young Professionals y Coordinación de actividades empresariales}

\vspace{0.3cm}

% Línea decorativa
\noindent\rule{2cm}{0.5pt}

\vfill

% Información inferior
{\color{sectionTeal}\bfseries IEEE Sección Aguascalientes}\\[0.2cm]
Última fecha de revisión: 6 de octubre 2026\\[0.1cm]
Fecha de autorización:

\newpage

% ============================================================
% DATOS DE IDENTIFICACIÓN
% ============================================================
\pagestyle{fancy}
\fancyhf{}
\renewcommand{\headrulewidth}{0pt}
\fancyfoot[C]{\thepage}

\vspace{0.5cm}

\seccionIEEE{Datos de identificación del documento}

\vspace{0.8cm}

\hyphenpenalty=10000
\exhyphenpenalty=10000

\begin{center}
\begin{tabular}{@{}p{4.5cm} p{8cm}@{}}
    \small Categoría & \small Detalle \\[0.5cm]
    
    \textbf{Título del Plan} & Fortalecimiento Comunitario y Sostenibilidad Operativa 2026 \\[0.5cm]
    
    \textbf{Sección IEEE} & Aguascalientes (Región 9) \\[0.5cm]
    
    \textbf{Coordinación Responsable} & Young Professionals y Actividades Empresariales \\[0.5cm]
    
    \textbf{Líder del Proyecto} & Ing. Tania Vanessa Ramírez Aguilar \\[0.5cm]
    
    \textbf{Fecha de Redacción} & Enero de 2026 \\[0.5cm]
    
    \textbf{Periodo de Aplicación} & Febrero - Febrero 2026 \\[0.5cm]
\end{tabular}
\end{center}


% ============================================================
% CONTENIDO
% ============================================================
\pagestyle{fancy}
\fancyhf{}
\renewcommand{\headrulewidth}{0pt}
\fancyfoot[C]{\thepage}

\vspace{0.5cm}

\section{Introducción y alcance estratégico}
El presente documento formaliza la estrategia de trabajo para el año 2026, impulsada por la coordinación de Young Professionals y la Coordinación de Actividades Empresariales de la Sección IEEE Aguascalientes.

El Plan de Trabajo se enfoca en la reactivación y el robustecimiento de la Sección, abordando la necesidad crítica de transformar la membresía en una comunidad de valor. Nuestro enfoque trasciende lo puramente técnico para concentrarse en la creación de espacios de convivencia óptimos y oportunidades de desarrollo profesional integral.

\textbf{Alcance Estratégico (Replicabilidad y Colaboración):}

Este plan no solo define las actividades operativas de 2026, sino que también establece una estructura metodológica y de gestión estandarizada para ser fácilmente replicada y transferida a futuros liderazgos. Además, la ejecución de las acciones 


\section{Objetivo principal}
El objetivo primordial de los planes de trabajo desarrollados es fomentar y fortalecer la comunidad entre los miembros, tanto de nuevo ingreso como los de trayectoria, con el fin estratégico de robustecer la Sección Aguascalientes.



\section{Justificación estratégica}
Este Plan de Trabajo se justifica en dos pilares fundamentales para el crecimiento sostenido de la Sección IEEE Aguascalientes: el desarrollo integral del capital humano y la sostenibilidad institucional.
\begin{itemize}
    \item Visión Humana (El Bienestar como Motor de la Excelencia Técnica): Reconocemos que la ingeniería es el vínculo que nos une, pero la base de nuestro propósito y acción es el ser humano integral. Las problemáticas identificadas por los jóvenes profesionales y estudiantes señalan una necesidad crítica de espacios de convivencia y networking óptimos. Fortalecer la comunidad es un requisito estratégico: un miembro satisfecho en sus necesidades de interacción y pertenencia es un profesional más motivado, productivo y comprometido con la excelencia técnica que promueve la IEEE.
\item Sostenibilidad y Transferencia de Conocimiento: Para asegurar el impacto a largo plazo, este documento se concibe con una estructura estandarizada y replicable. Se busca institucionalizar una metodología de trabajo que sea fácilmente transferible a futuras coordinaciones, garantizando la continuidad de las iniciativas de la Sección y la optimización del tiempo de los líderes entrantes.
\end{itemize}

\section{Líneas de acción y actividades}

La estrategia operativa se centra en tres líneas de acción fundamentales para el fortalecimiento de la comunidad.

\begin{itemize}[leftmargin=1.5cm]
    \item \textbf{Nota Estratégica:} La ejecución exitosa de estas actividades tiene un doble propósito: la autofinanciación de eventos (ej. Posada) y el aumento de la base de miembros para incrementar la asignación de fondos por parte de la IEEE a la Sección Aguascalientes.
\end{itemize}

\lineaAccion{I}{Reclutamiento}

\begin{tabularx}{\textwidth}{@{}p{3cm}X@{}}
    \filaDescriptiva{Evento}{Sesiones de difusión (2da Edición)-}
    
    \filaDescriptiva{Objetivo Principal}{\textbf{Reducir la brecha de acceso laboral,} garantizando a los jóvenes una interacción mínima de 20 minutos con representantes de empresas para generar un vínculo inicial, sin depender de recomendaciones.}
    
    \filaDescriptiva{Plan De Acción}{\textbf{Se pretende hacer un} listado de las carreras viables. Así mismo, ubicar profesores o compañeros conocidos de las mismas para facilitar esas reuniones.\newline\newline Modificar la presentación y perfeccionar.}
    
    \filaDescriptiva{Antecedente}{La edición anterior superó las expectativas de asistencia, demostrando la alta demanda por este formato.}
    
    \filaDescriptiva{Puntos de Mejora}{\textbf{Garantizar la puntualidad} estricta de las empresas.\newline\textbf{Solicitar un \textit{background} más específico} sobre los participantes.}
\end{tabularx}

\vspace{1cm}

\lineaAccion{II}{Capacitación}

\begin{tabularx}{\textwidth}{@{}p{3cm}X@{}}
    \filaDescriptiva{Evento}{Talleres técnicos y soft skills}
    
    \filaDescriptiva{Objetivo Principal}{\textbf{Desarrollar competencias profesionales} en los miembros mediante sesiones prácticas impartidas por expertos de la industria.}
    
    \filaDescriptiva{Plan De Acción}{Identificar las habilidades más demandadas en el mercado laboral actual.\newline\newline Establecer alianzas con empresas para obtener ponentes especializados.}
    
    \filaDescriptiva{Antecedente}{Los talleres anteriores tuvieron alta participación y retroalimentación positiva.}
    
    \filaDescriptiva{Puntos de Mejora}{\textbf{Diversificar los temas} para cubrir más áreas de interés.\newline\textbf{Mejorar la difusión} anticipada de los eventos.}
\end{tabularx}

\vspace{1cm}

\lineaAccion{III}{Networking}

\begin{tabularx}{\textwidth}{@{}p{3cm}X@{}}
    \filaDescriptiva{Evento}{Encuentros profesionales y sociales}
    
    \filaDescriptiva{Objetivo Principal}{\textbf{Fortalecer la comunidad IEEE} mediante eventos que fomenten la interacción entre miembros, egresados y profesionales de la industria.}
    
    \filaDescriptiva{Plan De Acción}{Organizar al menos un evento de networking por semestre.\newline\newline Invitar a egresados destacados como mentores.}
    
    \filaDescriptiva{Antecedente}{La comunidad ha expresado interés en más oportunidades de conexión profesional.}
    
    \filaDescriptiva{Puntos de Mejora}{\textbf{Establecer seguimiento} posterior a los eventos.\newline\textbf{Crear una base de datos} de contactos profesionales.}
\end{tabularx}

\end{document}
